\documentclass[11pt,preprint]{elsarticle}

\usepackage{lmodern}
%%%% My spacing
\usepackage{setspace}
\setstretch{1.2}
\DeclareMathSizes{12}{14}{10}{10}

% Wrap around which gives all figures included the [H] command, or places it "here". This can be tedious to code in Rmarkdown.
\usepackage{float}
\let\origfigure\figure
\let\endorigfigure\endfigure
\renewenvironment{figure}[1][2] {
    \expandafter\origfigure\expandafter[H]
} {
    \endorigfigure
}

\let\origtable\table
\let\endorigtable\endtable
\renewenvironment{table}[1][2] {
    \expandafter\origtable\expandafter[H]
} {
    \endorigtable
}


\usepackage{ifxetex,ifluatex}
\usepackage{fixltx2e} % provides \textsubscript
\ifnum 0\ifxetex 1\fi\ifluatex 1\fi=0 % if pdftex
  \usepackage[T1]{fontenc}
  \usepackage[utf8]{inputenc}
\else % if luatex or xelatex
  \ifxetex
    \usepackage{mathspec}
    \usepackage{xltxtra,xunicode}
  \else
    \usepackage{fontspec}
  \fi
  \defaultfontfeatures{Mapping=tex-text,Scale=MatchLowercase}
  \newcommand{\euro}{€}
\fi

\usepackage{amssymb, amsmath, amsthm, amsfonts}

\def\bibsection{\section*{References}} %%% Make "References" appear before bibliography


\usepackage[numbers]{natbib}

\usepackage{longtable}
\usepackage[margin=2.3cm,bottom=2cm,top=2.5cm, includefoot]{geometry}
\usepackage{fancyhdr}
\usepackage[bottom, hang, flushmargin]{footmisc}
\usepackage{graphicx}
\numberwithin{equation}{section}
\numberwithin{figure}{section}
\numberwithin{table}{section}
\setlength{\parindent}{0cm}
\setlength{\parskip}{1.3ex plus 0.5ex minus 0.3ex}
\usepackage{textcomp}
\renewcommand{\headrulewidth}{0.2pt}
\renewcommand{\footrulewidth}{0.3pt}

\usepackage{array}
\newcolumntype{x}[1]{>{\centering\arraybackslash\hspace{0pt}}p{#1}}

%%%%  Remove the "preprint submitted to" part. Don't worry about this either, it just looks better without it:
\makeatletter
\def\ps@pprintTitle{%
  \let\@oddhead\@empty
  \let\@evenhead\@empty
  \let\@oddfoot\@empty
  \let\@evenfoot\@oddfoot
}
\makeatother

 \def\tightlist{} % This allows for subbullets!

\usepackage{hyperref}
\hypersetup{breaklinks=true,
            bookmarks=true,
            colorlinks=true,
            citecolor=blue,
            urlcolor=blue,
            linkcolor=blue,
            pdfborder={0 0 0}}


% The following packages allow huxtable to work:
\usepackage{siunitx}
\usepackage{multirow}
\usepackage{hhline}
\usepackage{calc}
\usepackage{tabularx}
\usepackage{booktabs}
\usepackage{caption}


\newenvironment{columns}[1][]{}{}

\newenvironment{column}[1]{\begin{minipage}{#1}\ignorespaces}{%
\end{minipage}
\ifhmode\unskip\fi
\aftergroup\useignorespacesandallpars}

\def\useignorespacesandallpars#1\ignorespaces\fi{%
#1\fi\ignorespacesandallpars}

\makeatletter
\def\ignorespacesandallpars{%
  \@ifnextchar\par
    {\expandafter\ignorespacesandallpars\@gobble}%
    {}%
}
\makeatother


% definitions for citeproc citations
\NewDocumentCommand\citeproctext{}{}
\NewDocumentCommand\citeproc{mm}{%
\href{\#cite.\detokenize{#1}}{#2}\nocite{#1}}

\makeatletter
% allow citations to break across lines
\let\@cite@ofmt\@firstofone
% avoid brackets around text for \cite:
\def\@biblabel#1{}
\def\@cite#1#2{{#1\if@tempswa , #2\fi}}
\makeatother
\newlength{\cslhangindent}
\setlength{\cslhangindent}{1.5em}
\newlength{\csllabelwidth}
\setlength{\csllabelwidth}{3em}
\newenvironment{CSLReferences}[2] % #1 hanging-indent, #2 entry-spacing
{\begin{list}{}{%
	\setlength{\itemindent}{0pt}
	\setlength{\leftmargin}{0pt}
	\setlength{\parsep}{0pt}
	% turn on hanging indent if param 1 is 1
	\ifodd #1
	\setlength{\leftmargin}{\cslhangindent}
	\setlength{\itemindent}{-1\cslhangindent}
	\fi
	% set entry spacing
	\setlength{\itemsep}{#2\baselineskip}}}
{\end{list}}

\usepackage{calc}
\newcommand{\CSLBlock}[1]{\hfill\break\parbox[t]{\linewidth}{\strut\ignorespaces#1\strut}}
\newcommand{\CSLLeftMargin}[1]{\parbox[t]{\csllabelwidth}{\strut#1\strut}}
\newcommand{\CSLRightInline}[1]{\parbox[t]{\linewidth - \csllabelwidth}{\strut#1\strut}}
\newcommand{\CSLIndent}[1]{\hspace{\cslhangindent}#1}


\urlstyle{same}  % don't use monospace font for urls
\setlength{\parindent}{0pt}
\setlength{\parskip}{6pt plus 2pt minus 1pt}
\setlength{\emergencystretch}{3em}  % prevent overfull lines
\setcounter{secnumdepth}{5}

%%% Use protect on footnotes to avoid problems with footnotes in titles
\let\rmarkdownfootnote\footnote%
\def\footnote{\protect\rmarkdownfootnote}
\IfFileExists{upquote.sty}{\usepackage{upquote}}{}

%%% Include extra packages specified by user

%%% Hard setting column skips for reports - this ensures greater consistency and control over the length settings in the document.
%% page layout
%% paragraphs
\setlength{\baselineskip}{12pt plus 0pt minus 0pt}
\setlength{\parskip}{12pt plus 0pt minus 0pt}
\setlength{\parindent}{0pt plus 0pt minus 0pt}
%% floats
\setlength{\floatsep}{12pt plus 0 pt minus 0pt}
\setlength{\textfloatsep}{20pt plus 0pt minus 0pt}
\setlength{\intextsep}{14pt plus 0pt minus 0pt}
\setlength{\dbltextfloatsep}{20pt plus 0pt minus 0pt}
\setlength{\dblfloatsep}{14pt plus 0pt minus 0pt}
%% maths
\setlength{\abovedisplayskip}{12pt plus 0pt minus 0pt}
\setlength{\belowdisplayskip}{12pt plus 0pt minus 0pt}
%% lists
\setlength{\topsep}{10pt plus 0pt minus 0pt}
\setlength{\partopsep}{3pt plus 0pt minus 0pt}
\setlength{\itemsep}{5pt plus 0pt minus 0pt}
\setlength{\labelsep}{8mm plus 0mm minus 0mm}
\setlength{\parsep}{\the\parskip}
\setlength{\listparindent}{\the\parindent}
%% verbatim
\setlength{\fboxsep}{5pt plus 0pt minus 0pt}



\begin{document}



\begin{frontmatter}  %

\title{Report Loans and credit}

% Set to FALSE if wanting to remove title (for submission)




\author[Add1]{Tatjana Ebert}
\ead{}








\vspace{1cm}





\vspace{0.5cm}

\end{frontmatter}

\setcounter{footnote}{0}



%________________________
% Header and Footers
%%%%%%%%%%%%%%%%%%%%%%%%%%%%%%%%%
\pagestyle{fancy}
\chead{}
\rhead{}
\lfoot{}
\rfoot{\footnotesize Page \thepage}
\lhead{}
%\rfoot{\footnotesize Page \thepage } % "e.g. Page 2"
\cfoot{}

%\setlength\headheight{30pt}
%%%%%%%%%%%%%%%%%%%%%%%%%%%%%%%%%
%________________________

\headsep 35pt % So that header does not go over title




\section{Introduction}\label{introduction}

This report analyses household loan default patterns using anonymised
Lending Club data. The objective is to identify key predictors of loan
default, evaluate Lending Club's grading system, assess whether Texas
differs from other US states, and test the Director's existing
assumptions about default risk.

The main outcome variable is based on \texttt{loan\_status}. Loans
classified as \texttt{Fully\ Paid} are treated as successful repayments,
while loans classified as \texttt{Charged\ Off} or \texttt{Default} are
treated as written off as a loss. Loans listed as \texttt{Current} are
excluded from the main default analysis because their final outcome is
not yet known.

\section{Data and Cleaning}\label{data-and-cleaning}

\begin{longtable}[]{@{}
  >{\raggedleft\arraybackslash}p{(\linewidth - 8\tabcolsep) * \real{0.1912}}
  >{\raggedleft\arraybackslash}p{(\linewidth - 8\tabcolsep) * \real{0.1912}}
  >{\raggedleft\arraybackslash}p{(\linewidth - 8\tabcolsep) * \real{0.2353}}
  >{\raggedleft\arraybackslash}p{(\linewidth - 8\tabcolsep) * \real{0.2647}}
  >{\raggedleft\arraybackslash}p{(\linewidth - 8\tabcolsep) * \real{0.1176}}@{}}
\caption{Summary statistics for cleaned Lending Club
data}\tabularnewline
\toprule\noalign{}
\begin{minipage}[b]{\linewidth}\raggedleft
observations
\end{minipage} & \begin{minipage}[b]{\linewidth}\raggedleft
default\_rate
\end{minipage} & \begin{minipage}[b]{\linewidth}\raggedleft
avg\_loan\_amount
\end{minipage} & \begin{minipage}[b]{\linewidth}\raggedleft
avg\_interest\_rate
\end{minipage} & \begin{minipage}[b]{\linewidth}\raggedleft
avg\_dti
\end{minipage} \\
\midrule\noalign{}
\endfirsthead
\toprule\noalign{}
\begin{minipage}[b]{\linewidth}\raggedleft
observations
\end{minipage} & \begin{minipage}[b]{\linewidth}\raggedleft
default\_rate
\end{minipage} & \begin{minipage}[b]{\linewidth}\raggedleft
avg\_loan\_amount
\end{minipage} & \begin{minipage}[b]{\linewidth}\raggedleft
avg\_interest\_rate
\end{minipage} & \begin{minipage}[b]{\linewidth}\raggedleft
avg\_dti
\end{minipage} \\
\midrule\noalign{}
\endhead
\bottomrule\noalign{}
\endlastfoot
383202 & 0.225 & 14544.93 & 13.024 & 18.792 \\
\end{longtable}

\section{Exploratory Analysis}\label{exploratory-analysis}

Default rates differ substantially across borrower characteristics. The
first important comparison is by credit grade.

\begin{longtable}[]{@{}lrr@{}}
\caption{Default rates by Lending Club credit grade}\tabularnewline
\toprule\noalign{}
grade & default\_rate & n \\
\midrule\noalign{}
\endfirsthead
\toprule\noalign{}
grade & default\_rate & n \\
\midrule\noalign{}
\endhead
\bottomrule\noalign{}
\endlastfoot
A & 0.069 & 66886 \\
B & 0.156 & 114785 \\
C & 0.253 & 111454 \\
D & 0.346 & 54883 \\
E & 0.431 & 24302 \\
F & 0.524 & 8567 \\
G & 0.573 & 2325 \\
\end{longtable}

\begin{figure}[H]
\centering
\includegraphics[width=0.8\textwidth]{C:/Users/Tatjana/Desktop/Stellenbosch/Uni/Kurse/Data Science/Data Part A/Exam/31074774/Question_3_Loans/graphs/Grade.png}
\caption{Default rates by credit grade}
\end{figure}

Debt-to-income ratios are also central to the analysis because they
measure the borrower's repayment burden relative to income.

\begin{longtable}[]{@{}lrr@{}}
\caption{Default rates by DTI band}\tabularnewline
\toprule\noalign{}
dti\_band & default\_rate & n \\
\midrule\noalign{}
\endfirsthead
\toprule\noalign{}
dti\_band & default\_rate & n \\
\midrule\noalign{}
\endhead
\bottomrule\noalign{}
\endlastfoot
{[}0,5{]} & 0.169 & 19605 \\
(5,10{]} & 0.167 & 48750 \\
(10,15{]} & 0.186 & 74676 \\
(15,20{]} & 0.212 & 79480 \\
(20,25{]} & 0.241 & 67092 \\
(25,30{]} & 0.279 & 49716 \\
(30,35{]} & 0.316 & 30001 \\
(35,40{]} & 0.321 & 10123 \\
(40,45{]} & 0.340 & 1573 \\
(45,50{]} & 0.356 & 947 \\
(50,55{]} & 0.248 & 238 \\
(55,60{]} & 0.226 & 199 \\
\end{longtable}

\begin{figure}[H]
\centering
\includegraphics[width=0.8\textwidth]{C:/Users/Tatjana/Desktop/Stellenbosch/Uni/Kurse/Data Science/Data Part A/Exam/31074774/Question_3_Loans/graphs/cap_dti.png}
\caption{Default rates by DTI cap}
\end{figure}

\section{Testing the Institute's
Heuristics}\label{testing-the-institutes-heuristics}

\subsection{Homeowners and Long
Employment}\label{homeowners-and-long-employment}

The first belief is that homeowners and borrowers employed for more than
ten years have a much lower default probability on short-term loans.
Short-term loans are defined as 36-month loans.

\begin{longtable}[]{@{}rrrr@{}}
\caption{Default rates for short-term loans by home ownership and
employment length}\tabularnewline
\toprule\noalign{}
homeowner & emp\_10plus & default\_rate & n \\
\midrule\noalign{}
\endfirsthead
\toprule\noalign{}
homeowner & emp\_10plus & default\_rate & n \\
\midrule\noalign{}
\endhead
\bottomrule\noalign{}
\endlastfoot
0 & 0 & 0.232 & 92813 \\
0 & 1 & 0.230 & 28042 \\
1 & 0 & 0.172 & 107966 \\
1 & 1 & 0.146 & 68636 \\
\end{longtable}

\begin{figure}[H]
\centering
\includegraphics[width=0.8\textwidth]{C:/Users/Tatjana/Desktop/Stellenbosch/Uni/Kurse/Data Science/Data Part A/Exam/31074774/Question_3_Loans/graphs/homeowners.png}
\caption{Short-term default rates by home ownership and employment length}
\end{figure}

```

This comparison shows whether the Institute's belief is supported in the
raw data. However, raw differences may be misleading because homeowners
and long-employed borrowers may also differ in credit grade, income and
DTI.

\subsection{State Differences}\label{state-differences}

The second belief is that US states differ significantly in their
culture of defaulting on short-term loans. The table below compares
state-level default rates among states with sufficient observations.

\begin{longtable}[]{@{}lrr@{}}
\caption{Highest short-term default rates by state}\tabularnewline
\toprule\noalign{}
addr\_state & default\_rate & n \\
\midrule\noalign{}
\endfirsthead
\toprule\noalign{}
addr\_state & default\_rate & n \\
\midrule\noalign{}
\endhead
\bottomrule\noalign{}
\endlastfoot
AR & 0.236 & 2198 \\
MS & 0.228 & 1885 \\
OK & 0.225 & 2748 \\
LA & 0.225 & 3263 \\
AL & 0.224 & 3579 \\
NE & 0.222 & 1470 \\
NY & 0.214 & 23664 \\
FL & 0.211 & 22551 \\
NV & 0.206 & 4803 \\
NJ & 0.202 & 10246 \\
HI & 0.201 & 1387 \\
AK & 0.199 & 643 \\
MO & 0.199 & 4617 \\
MD & 0.198 & 6540 \\
OH & 0.197 & 9675 \\
\end{longtable}

Texas is examined separately because the Director is based in Texas.

\begin{longtable}[]{@{}rrr@{}}
\caption{Default rate in Texas versus other states}\tabularnewline
\toprule\noalign{}
texas & default\_rate & n \\
\midrule\noalign{}
\endfirsthead
\toprule\noalign{}
texas & default\_rate & n \\
\midrule\noalign{}
\endhead
\bottomrule\noalign{}
\endlastfoot
0 & 0.224 & 350819 \\
1 & 0.228 & 32383 \\
\end{longtable}

\subsection{Age and Occupation}\label{age-and-occupation}

The dataset does not contain a direct borrower age variable. Therefore,
the claim that credit grades predict whether younger individuals default
cannot be tested directly. Similarly, if occupation is absent from the
dataset, interest-rate determination by occupation cannot be evaluated
directly. This is an important data limitation.

\section{Econometric Model}\label{econometric-model}

A logistic regression is appropriate because the dependent variable is
binary: default or no default.

\begin{longtable}[]{@{}lrrr@{}}
\caption{Logistic regression results for default
probability}\tabularnewline
\toprule\noalign{}
term & estimate & odds\_ratio & p.value \\
\midrule\noalign{}
\endfirsthead
\toprule\noalign{}
term & estimate & odds\_ratio & p.value \\
\midrule\noalign{}
\endhead
\bottomrule\noalign{}
\endlastfoot
(Intercept) & -3.287 & 0.037 & 0.000 \\
gradeB & 0.637 & 1.891 & 0.000 \\
gradeC & 1.011 & 2.749 & 0.000 \\
gradeD & 1.226 & 3.409 & 0.000 \\
gradeE & 1.340 & 3.818 & 0.000 \\
gradeF & 1.513 & 4.541 & 0.000 \\
gradeG & 1.573 & 4.820 & 0.000 \\
term60 months & 0.407 & 1.503 & 0.000 \\
int\_rate & 0.033 & 1.034 & 0.000 \\
dti & 0.008 & 1.008 & 0.000 \\
annual\_inc & 0.000 & 1.000 & 0.000 \\
home\_ownershipMORTGAGE & -0.205 & 0.814 & 0.483 \\
home\_ownershipOWN & 0.008 & 1.008 & 0.977 \\
home\_ownershipRENT & 0.252 & 1.286 & 0.390 \\
emp\_length1 year & 0.031 & 1.031 & 0.146 \\
emp\_length10+ years & -0.044 & 0.957 & 0.007 \\
emp\_length2 years & -0.005 & 0.995 & 0.808 \\
emp\_length3 years & 0.017 & 1.018 & 0.392 \\
emp\_length4 years & -0.001 & 0.999 & 0.974 \\
emp\_length5 years & 0.011 & 1.011 & 0.629 \\
emp\_length6 years & -0.046 & 0.955 & 0.061 \\
emp\_length7 years & -0.015 & 0.986 & 0.585 \\
emp\_length8 years & 0.011 & 1.011 & 0.651 \\
emp\_length9 years & 0.007 & 1.007 & 0.781 \\
emp\_lengthn/a & 0.462 & 1.587 & 0.000 \\
purposecredit\_card & 0.205 & 1.227 & 0.000 \\
purposedebt\_consolidation & 0.241 & 1.273 & 0.000 \\
purposehome\_improvement & 0.272 & 1.313 & 0.000 \\
purposehouse & 0.090 & 1.094 & 0.177 \\
purposemajor\_purchase & 0.259 & 1.296 & 0.000 \\
\end{longtable}

The model estimates how borrower characteristics are associated with the
probability of default. Coefficients above zero increase default risk,
while coefficients below zero reduce default risk. Odds ratios above one
indicate higher default odds.

\section{Model Performance}\label{model-performance}

\begin{longtable}[]{@{}r@{}}
\caption{Model AUC}\tabularnewline
\toprule\noalign{}
AUC \\
\midrule\noalign{}
\endfirsthead
\toprule\noalign{}
AUC \\
\midrule\noalign{}
\endhead
\bottomrule\noalign{}
\endlastfoot
0.709 \\
\end{longtable}

The AUC measures the model's ability to distinguish between defaulted
and fully paid loans. An AUC of 0.5 indicates no predictive power, while
values closer to 1 indicate stronger classification performance.

\section{DTI Hard-Cap Recommendation}\label{dti-hard-cap-recommendation}

A hard cap on DTI can reduce risk by excluding borrowers whose debt
burden is too high relative to income. The appropriate cap depends on
the lender's default tolerance.

\begin{longtable}[]{@{}lrr@{}}
\caption{DTI bands and associated default rates}\tabularnewline
\toprule\noalign{}
dti\_band & default\_rate & n \\
\midrule\noalign{}
\endfirsthead
\toprule\noalign{}
dti\_band & default\_rate & n \\
\midrule\noalign{}
\endhead
\bottomrule\noalign{}
\endlastfoot
{[}0,5{]} & 0.169 & 19605 \\
(5,10{]} & 0.167 & 48750 \\
(10,15{]} & 0.186 & 74676 \\
(15,20{]} & 0.212 & 79480 \\
(20,25{]} & 0.241 & 67092 \\
(25,30{]} & 0.279 & 49716 \\
(30,35{]} & 0.316 & 30001 \\
(35,40{]} & 0.321 & 10123 \\
(40,45{]} & 0.340 & 1573 \\
(45,50{]} & 0.356 & 947 \\
(50,55{]} & 0.248 & 238 \\
(55,60{]} & 0.226 & 199 \\
\end{longtable}

Based on this table, a conservative lender should choose a lower DTI
cap, while a lender with greater risk tolerance may accept higher DTI
levels. A reasonable policy structure is:

\begin{itemize}
\tightlist
\item
  Low default tolerance: DTI cap around 20 to 25.
\item
  Moderate default tolerance: DTI cap around 30.
\item
  High default tolerance: DTI cap around 35 to 40.
\end{itemize}

The exact recommendation should be based on where the observed default
rate begins to rise sharply in the DTI table and graph.

\section{Texas Compared with Other
States}\label{texas-compared-with-other-states}

\begin{longtable}[]{@{}lrrrrr@{}}
\caption{Controlled regression test for Texas effect}\tabularnewline
\toprule\noalign{}
term & estimate & std.error & statistic & p.value & odds\_ratio \\
\midrule\noalign{}
\endfirsthead
\toprule\noalign{}
term & estimate & std.error & statistic & p.value & odds\_ratio \\
\midrule\noalign{}
\endhead
\bottomrule\noalign{}
\endlastfoot
(Intercept) & -2.825 & 0.027 & -104.143 & 0.000 & 0.059 \\
texas & 0.033 & 0.014 & 2.282 & 0.022 & 1.034 \\
gradeB & 0.742 & 0.020 & 37.646 & 0.000 & 2.099 \\
gradeC & 1.163 & 0.025 & 45.656 & 0.000 & 3.200 \\
gradeD & 1.406 & 0.036 & 38.848 & 0.000 & 4.080 \\
gradeE & 1.550 & 0.047 & 32.927 & 0.000 & 4.709 \\
gradeF & 1.746 & 0.059 & 29.696 & 0.000 & 5.731 \\
gradeG & 1.817 & 0.076 & 24.026 & 0.000 & 6.151 \\
dti & 0.007 & 0.000 & 18.449 & 0.000 & 1.008 \\
int\_rate & 0.035 & 0.003 & 12.621 & 0.000 & 1.036 \\
term60 months & 0.292 & 0.010 & 30.014 & 0.000 & 1.339 \\
annual\_inc & 0.000 & 0.000 & -19.307 & 0.000 & 1.000 \\
\end{longtable}

If the Texas coefficient is statistically significant after controlling
for borrower risk characteristics, then Texas appears different from
other states. If it is not significant, then any raw Texas difference is
likely explained by borrower composition rather than a distinct default
culture.

\section{Discussion}\label{discussion}

The key risk drivers expected from this analysis are credit grade,
interest rate, DTI, loan term, annual income, delinquencies and recent
credit inquiries. Borrowers with weaker credit grades, higher DTI,
higher interest rates and longer loan terms are expected to show higher
default probabilities.

Lending Club's grading system is useful if default rates increase
consistently from grade A to lower grades. However, the grading system
should not be used alone. It should be combined with DTI, income, loan
purpose, term and prior credit behaviour.

\section{Conclusion}\label{conclusion}

The analysis provides evidence on which borrower characteristics are
most strongly associated with loan default. The most relevant risk
indicators are expected to be credit grade, DTI, interest rate and loan
term. A DTI cap can help reduce default risk, but the appropriate
threshold depends on the lender's tolerance for defaults. Texas should
be evaluated using both raw default rates and controlled regression
results.

The Institute's heuristics are only partly testable. The homeowner and
employment-length claim can be tested among short-term loans. State
differences can be examined both descriptively and with controls.
However, the claim about younger borrowers cannot be tested directly
because the dataset does not contain borrower age.

\bibliography{Tex/ref}





\end{document}
